% !PRINTED 248
\weberpage{248}
Für die Funktion $\gamma_3(\omega)$ ist [nach § 54, (15)]:
\[
\gamma_3\left(\frac{\gamma+\delta\omega}{\alpha+\beta\omega}\right) = (-1)^{\alpha\gamma+\beta\gamma+\beta\delta} \gamma_3(\omega),
\]
und in den Zusammensetzungen (8), (10), § 69 ist
\[
\lambda \equiv 1, \quad \alpha \equiv \delta \equiv 0 \ (\mathrm{mod}\ 2).
\]
Die $\nu$ Größen
\weq{2}{\gamma_3\left(\frac{c+\partial\omega}{a}\right) \gamma_3(\omega)}
vertauschen sich daher untereinander und sind also gleichfalls die
Wurzeln einer Transformationsgleichung erster Stufe.

Um für den einfachsten Fall $n=2$ die erstere dieser Gleichungen zu
bilden, beachte man die Relation [§ 54, (5)]:
\weq{3}{\gamma_2(\omega) = \frac{f(\omega)^{24}-16}{f(\omega)^8} = \frac{f_1(\omega)^{24}+16}{f_1(\omega)^8} = \frac{f_2(\omega)^{24}+16}{f_2(\omega)^8},}
woraus wegen
\[
f_1(2\omega) f_2(\omega) = \sqrt{2} \qquad [\S\ 34, (16)]
\]
folgt:
\weq{4}{
\begin{gathered}
\gamma_2(2\omega) = \frac{2^8+f_2(\omega)^{24}}{f_2(\omega)^{16}} \\[4pt]
\gamma_2\left(\frac{\omega}{2}\right) = \frac{2^8+f_1(\omega)^{24}}{f_1(\omega)^{16}} \\[4pt]
\gamma_2\left(\frac{\omega+3}{2}\right) = \frac{2^8-f(\omega)^{24}}{f(\omega)^{16}}. \qquad [\S\ 34, (13)]
\end{gathered}
}

Bezeichnen wir diese drei Größen mit $x$, $x_0$, $x_1$, so ergibt
sich aus den Relationen (6), (8), § 54:
\[
\begin{aligned}
x + x_0 + x_1 &= \gamma_2(\omega)^2 \\
xx_0 + xx_1 + x_0x_1 &= 495\,\gamma_2(\omega) \\
xx_0x_1 &= -j(\omega) + 2^4\cdot3^3\cdot5^3,
\end{aligned}
\]
so daß $x$, $x_0$, $x_1$ die Wurzeln der Gleichung sind
\weq{5}{x^3 - \gamma_2(\omega)^2 x^2 + 5\cdot9\cdot11\cdot\gamma_2(\omega) x + j(\omega) - 2^4\cdot3^3\cdot5^3 = 0.}
Die Gleichung, deren Wurzeln die Kuben der Wurzeln von (5) sind, ist
die Invariantengleichung für $n=2$, und läßt sich daraus ohne
Schwierigkeit berechnen.

\begin{center}
\textbf{§ 72. Multiplikatorgleichungen erster Stufe.}
\end{center}

Unter den Multiplikatorgleichungen sollen hier die betrachtet werden,
deren Wurzeln die verschiedenen Potenzen von $P_{11}$ sind,
multipliziert mit Potenzen von $\gamma_2(\omega)$, $\gamma_3(\omega)$,
deren Koeffizienten
%CONT>>
